% Part: first-order-logic
% Chapter: model-theory
% Section: isomorphism

\documentclass[../../../include/open-logic-section]{subfiles}

\begin{document}

\olfileid{mod}{bas}{iso}
\olsection{同构结构}

一阶结构有两种彼此相似的方式。一种相似性在于它们满足相同的语句。我们称这样的结构是\emph{初等等价的}。但结构可以差异巨大，却仍然满足相同的语句——例如，一个结构可数而另一个不可数。这是因为结构的许多特征无法在一阶语言中表达，原因要么在于语言不够丰富，要么在于一阶逻辑的根本局限性（如 勒文海姆--斯科伦 定理）。因此，结构之间另一种更严格的相似性在于它们本质上是相同的，即它们仅在构成元素上不同，而在结构特征上完全一致。将其精确化的方法是借助\emph{同构}的概念。

\begin{defn}
\ollabel{defn:elem-equiv}
给定同一语言~$\Lang{L}$ 的两个结构 $\Struct{M}$ 和 $\Struct M'$，我们称 $\Struct{M}$ \emph{初等等价于} $\Struct M'$，记作 $\Struct{M} \equiv \Struct M'$，当且仅当对于 $\Lang{L}$ 的每个语句~$!A$，$\Sat{M}{!A}$ 当且仅当 $\Sat{M'}{!A}$。
\end{defn}

\begin{defn}
\ollabel{defn:isomorphism}
给定同一语言~$\Lang L$ 的两个结构 $\Struct{M}$ 和 $\Struct M'$，我们称 $\Struct{M}$ \emph{同构于} $\Struct M'$，记作 $\Struct{M} \simeq \Struct M'$，当且仅当存在一个函数 $h \colon
\Domain{M} \to \Domain{M'}$ 满足：
\begin{enumerate}
\item $h$ 是单射的：如果 $h(x) = h(y)$，那么 $x = y$； 
\item $h$ 是满射的：对于每个 $y \in \Domain{M'}$，都存在 $x \in \Domain{M}$ 使得 $h(x) = y$；
\item \ollabel{defn:iso-const}对于每个常项符号 $c$：$h(\Assign{c}{M}) = \Assign{c}{M'}$；
\item \ollabel{defn:iso-pred}对于每个 $n$ 元谓词符号~$P$：\[
  \tuple{a_1, \dots, a_n}\in \Assign{P}{M} \quad\text{当且仅当}\quad
  \tuple{h(a_1), \dots, h(a_n)} \in \Assign{P}{M'};
  \]
\item \ollabel{defn:iso-func}对于每个 $n$ 元函数符号 $f$：\[
  h(\Assign{f}{M}(a_1, \dots, a_n)) =
  \Assign{f}{M'}(h(a_1), \dots, h(a_n)).
  \]
\end{enumerate}
\end{defn}

\begin{thm}
\ollabel{thm:isom}
如果 $\Struct{M} \iso \Struct M'$，那么 $\Struct{M} \elemequiv \Struct{M'}$。
\end{thm}

\begin{proof}
设 $h$ 是 $\Struct{M}$ 到 $\Struct M'$ 的一个同构。对于任意指派~$s$，$h \circ s$ 是 $h$ 和 $s$ 的复合，即 $\Struct{M'}$ 中满足 $(h \circ s)(x) = h(s(x))$ 的指派。通过对 $t$ 和 $!A$ 归纳，可以证明以下更强的断言：
\begin{enumerate}
  \item[a.] $h(\Value{t}{M}[s]) = \Value{t}{M'}[h\circ s]$。
  \item[b.] $\Sat{M}{!A}[s]$ 当且仅当 $\Sat{M'}{!A}[h \circ s]$。
\end{enumerate}
第一条通过对 $t$ 的复杂度进行归纳来证明。
\begin{enumerate}
\item 如果 $t \ident c$，那么 $\Value{c}{M}[s] = \Assign{c}{M}$ 且 $\Value{c}{M'}[h \circ s] = \Assign{c}{M'}$。因此，$h(\Value{t}{M}[s]) = h(\Assign{c}{M}) = \Assign{c}{M'}$（据 \olref{defn:isomorphism} 的 \olref{defn:iso-const}）$=
  \Value{t}{M'}[h \circ s]$。
\item 如果 $t \ident x$，那么 $\Value{x}{M}[s] = s(x)$ 且 $\Value{x}{M'}[h \circ s] = h(s(x))$。因此，$h(\Value{x}{M}[s]) =
  h(s(x)) = \Value{x}{M'}[h \circ s]$。
\item 如果 $t \ident f(t_1, \dots, t_n)$，那么 \begin{align*}
    \Value{t}{M}[s] & = \Assign{f}{M}(\Value{t_1}{M}[s], \dots, \Value{t_n}{M}[s]) \quad\text{且}\\
  \Value{t}{M'}[h \circ s] & = \Assign{f}{M'}(\Value{t_1}{M'}[h \circ
    s], \dots, \Value{t_n}{M'}[h \circ s]).
  \end{align*}
  归纳假设是对于每个 $i$，$h(\Value{t_i}{M}[s])
  = \Value{t_i}{M'}[h\circ s]$。所以，\begin{align}
    h(\Value{t}{M}[s]) 
    & = h(\Assign{f}{M}(\Value{t_1}{M}[s], \dots, \Value{t_n}{M}[s]))\notag\\
    & = \Assign{f}{M'}(h(\Value{t_1}{M}[s]), \dots,
    h(\Value{t_n}{M}[s])) \ollabel{iso-1}\\
    & = \Assign{f}{M'}(\Value{t_1}{M'}[h \circ s], \dots,
    \Value{t_n}{M'}[h \circ s]) \ollabel{iso-2}\\
    & = \Value{t}{M'}[h\circ s] \notag
  \end{align}
  这里，\olref{iso-1} 由 \olref{defn:isomorphism} 的 \olref{defn:iso-func} 得出，而 \olref{iso-2} 由归纳假设得出。
\end{enumerate}
(b) 部分的证明留作练习。

如果 $!A$ 是一个语句，那么指派~$s$ 和 $h \circ s$ 就不再重要，从而 $\Sat{M}{!A}$ 当且仅当 $\Sat{M'}{!A}$。
\end{proof}

\begin{prob}
详细完成 \olref[mod][bas][iso]{thm:isom} 中 (b) 部分的证明。请务必指出 \olref[mod][bas][iso]{defn:isomorphism} 中刻画同构的五条性质分别在何处被用到。
\end{prob}

\begin{defn}
结构 $\Struct{M}$ 的一个\emph{自同构}是 $\Struct{M}$ 到其自身的一个同构。
\end{defn}

\begin{prob}
证明对于任意结构 $\Struct{M}$，如果 $X$ 是 $\Struct{M}$ 的一个可定义子集，且 $h$ 是 $\Struct{M}$ 的一个自同构，那么 $X
= \Setabs{h(x)}{x \in X}$（即 $X$ 在 $h$ 下保持不变）。
\end{prob}


\end{document}